\section{The Seiberg-Witten Analysis of $N=2$ Gauge Theory With Matter}

In this section, we consider the $N=2$ supersymmetric gauge theory,
with gauge group $SU(2)$, coupled to $N_f$ matter multiplets. The
analysis involves many technical issues and the properties of the
theory depend on the value of $N_f$. As in usual QCD, matter fields
(fermions) and gauge bosons contribute to the $\beta$-function with
opposite signs. Thus, to insure asymptotic freedom, we require $N_f
\leq 4$. The theory with $N_f =4$ is finite to all orders in
perturbation theory and, when the quarks are massless, it is very
likely conformally invariant non-perturbatively. Seiberg and Witten
provided evidence that this theory is an example of an $SL(2,Z)$
invariant theory. The only other $SL(2,Z)$ invariant theory known in
four dimensions is the $N=4$ supersymmetric gauge theory which also
has a vanishing $\beta$-function. However, in the following, we will
mainly concentrate on $N_f\leq 3$. For the $N_f=4$ case, the reader is
referred to original work of Seiberg and Witten \cite{SW2}.

In the presence of matter, it is convenient to choose a charge
normalization different from the previous sections. In $N=2$ pure
gauge theory all fields transform in the adjoint representation and
the charges of particles and monopoles are integers. Hence, in the
formula $Z=a n_e + a_D n_m$, $n_e$ and $n_m$ are integers. With this
normalization, when quarks are present, $n_e$ could be half-integral.
In this section we choose a slightly different convention: To ensure
that $n_e$ is always integral, we multiply it by 2 and divide $a$ by 2
to keep the mass formula unchanged. Since $a_D$ is kept unchanged, it
is now given by $2 a_D=\partial{\cal F}/\partial a$. In terms of the
gauge invariant quantity $u={\rm Tr}\phi^2$, as $|u|\rightarrow\infty$,
the asymptotic behaviour of $a$ and $a_D$ with this normalization
becomes 
$$
a \cong \frac{1}{2} \sqrt{2u}\,, \quad
a_D \simeq i \frac{4}{\pi} a \log a\,.
$$
Furthermore, the effective coupling is also rescaled to 
$\tau=\partial a_D/\partial a = \frac{\theta}{\pi} + \frac{8 \pi
i}{g^2}$. This rescaling of the charge affects the form of the curve
which determines the solution of the model, keeping the physics
unchanged.  

In this section we are going to study the theory with gauge group
$SU(2)$. In this case, the most interesting dynamics appear in the
Coulomb branch of the moduli space, where $SU(2)$ breaks down to 
$U(1)$. It is, therefore, instructive to first study $N = 2$
supersymmetric QED.

\subsection{$N = 2$ Supersymmetric QED}

The $N=2$ supersymmetric QED contains the photon vector multiplet and
$k$ quark hypermultiplets. The photon multiplet consists of the $N=1$
chiral superfields $(W_{\alpha}, A)$ and each quark hypermultiplet
consists of two $N=1$ chiral superfields $M^i$ and $\widetilde M_i$,
with $U(1)$ charges $+1$ and $-1$, respectively. The only
renormalizable $N = 2$ compatible superpotential is: 
$$
{\cal W}=\sqrt{2} A M^i\widetilde M_i + \sum_i m_i M^i\widetilde
M_i\,. 
$$
When $m_i = 0$, the theory has a global symmetry group $SU(k)\times
SU(2)_R \times U(1)_{\cal R}$. $SU(k)$ is a flavour symmetry group
acting on the $k$ hypermultiplets, while $SU(2)_R$ rotates the two
supersymmetries and $U(1)_{\cal R}$ is the usual $R$-symmetry which is
afflicted by an anomaly. As in the case without matter, the strategy
is to first determine the classical moduli space and then to
see whether the quantum moduli space is qualitatively different.
The symmetries of the theory are very useful in determining the
structure of these moduli spaces.

In general, the classical moduli space has a Higgs branch and a
Coulomb branch. The Higgs branch is defined by some of the $M^i$'s
acquiring vacuum expectation values. In this case, both in QED and in
QCD with gauge group $SU(2)$, the gauge group is completely broken.
For $k=1$, there is no Higgs phase.  When $k \geq 2$, and $m_i = 0$,
then, using the global symmetries together with the the vanishing
condition for the $D$-term, it is possible to rotate $M$ and
$\widetilde M$ to the form $M=(B,0,0,...)$ and $\widetilde M = (0,
B,0,...)$. From this one can read off the different patterns of
symmetry breaking as a function of the number of flavours $k$. A
theory on the Higgs branch does not contain monopoles or dyons and
hence, the dynamical possibilities are not as rich as on the Coulomb
branch. The classical moduli space is a hyperK\"ahler manifold and the
symmetries of the theory lead to a unique hyperK\"ahler metric on
it. As a result of this uniqueness, the metric does not receive
quantum mechanical corrections. Therefore, on the Higgs branch, the
classical and the quantum moduli spaces are the same.

In the case at hand, if all $m_i = 0$, then the Coulomb branch is
defined by $<A> \not= 0$ which, in turn, implies that
$<M^i>=<\widetilde M_i>=0$. The $U(1)$ gauge group remains unbroken
while all the $M$'s become massive. At a generic point on the moduli
space, the effective low-energy threory (which involves only the
massless modes) is a pure $N=2$ gauge theory and the K\"ahler
potential is of the special geometry type : $K={\rm Im} a_D (a)
\bar{a}$. This is related by $N=2$ supersymmetry to the the gauge
kinetic term,  
$$
\int d^2 \theta \; \frac{\partial a_D}{\partial a} \;
W^{\alpha}W_{\alpha}\,.
$$
The K\"ahler potential (and thus the metric) receive quantum
corrections, but the one-loop approximation to $K$ is exact and also
there are no non-perturbative corrections since this theory does not
contain instantons. Using the one-loop $\beta$-function for QED with
$k$ hypermultiplets, we obtain 
\begin{equation}
a_D = - \frac{ik}{2\pi} \; a \ln (a/\Lambda)\,.
\label{beta-hyper}
\end{equation}
To reproduce this formula one should be careful about the extra
factors of $2$ in the supersymmetric way of defining $g$ and the
normalization convention mentioned earlier. The metric ${\rm Im}
(\tau)$ obtained from this is zero at $|a| = \Lambda / e$ and thus the
effective coupling constant is singular. Because of this Landau pole
singularity, the theory does not make sense in the ultra-violet
region unless embedded in a larger theory which is asymptotically
free. If fermion mass terms are added, the singularities on the  
Coulomb branch can move.  Since, when $a = - \frac{1}{\sqrt{2}} m_i$
one electron becomes massless, we have
$$
a_D = - \frac{i}{2\pi} \sum_i (a+m_i / \sqrt{2}) \;\ln
\frac{(a+m_i / \sqrt{2})}{\Lambda}\,.
$$
To each bare mass is now associated a singularity on the moduli space
where the particle becomes massless.  Depending on the possible
equality of two or more masses, one can have Higgs and Coulomb
branches touching, leading to an intricate structure.

Here, we also see a manifestation of the modified form of the central
charge of the $N = 2$ algebra discussed in section 2. For pure $N=2$,
we have $Z = n_e a + n_m a_D$. However, now the electron masses are
not just $\sqrt{2} |a|$, rather the $i$-th multiplet has mass 
$|\sqrt{2} a + m_i |$. As shown in section 2, this is consistent with
the fact that the $U(1)$ charges $S_i$ of the massive hypermultiplets
appears in the formula for the central charge $Z$ as
$$
Z = n_e a + n_m a_D + \sum_i S_i m_i / \sqrt{2}\,.
$$

\subsection{$N=2$ Supersymmetric QCD with Matter}
The $N=2$ supersymmetric QCD coupled to $N_f$ matter hypermultiplets 
contains the $N=1$ superpotential
$$
W=\sum_{i=1}^{N_f}(\sqrt{2}\widetilde{Q}_i\Phi
Q^i+m_i\tilde{Q}_iQ^i)\,. 
$$
In general, when $m_i=0$, this theory has a global symmetry group
$SU(N_f)\times SU(2)_R\times U(1)_{\cal R}$. In the special case, when
the gauge group is $SU(2)$, the flavour group is enlarged to
$O(2N_f)$. This is due to the fact that for $SU(2)$ the fundamental
representation and its conjugate are isomorphic and, therefore, $Q^i$
and $\widetilde Q_i$ can be combined into a $2N_f$-dimensional vector
transforming under $O(2N_f)$. Thus, for the special case of the
$SU(2)$ gauge group, the theory also has a parity symmetry group ${\bf
Z}_2\subset O(2N_f)$ acting as
\begin{equation}
\rho : Q_1 \leftrightarrow \widetilde{Q}_1\,,
\label{rho}
\end{equation}
with all other fields remaining unchanged. This parity plays an
important role in the analysis of the theory.  The global symmetry
group of the theory is actually a quotient of $O(2N_f) \times SU(2)_R
\times U(1)_{\cal R}$. The quotient is to be taken because a $Z_2
\subset U(1)_{\cal R}$ is the same as $(-1)^F$ contained in the
Lorentz group. This, combined with the center of $SU(2)_R$, is the
same as the ${\bf Z}_2$ in the center of $O(2N_f)$.

For $N_f = 0,1$, the theory has only a Coulomb branch with $(<\phi>
\not= 0)$. On this branch $SU(2)$ is broken to $U(1)$ and all quarks
are massive. Moreover, $U(1)_{\cal R}$ is spontaneously broken because
$\Phi$ has $U(1)_{\cal R}$ charge $2$. For $N_f \geq 2$ we can either
have a Coulomb branch or Higgs  branches. On a Higgs branch, the
gauge symmetry is fully broken and the dynamics is not very rich. We
will not analyse this branch in detail, but will only mention that, as
in the QED case, the metric on the moduli space is uniquely
determined by the symmetries and does not receive quantum correction. 

\noindent\underline{Some Properties of the Quantum Theory}: 
The perturbative $\beta$-function of our theory (which, due to
supersymmetry, is only a one-loop effect) is given by
\begin{equation}
\beta (g) = - \frac{4-N_f}{16\pi^2} g^3.
\label{betanf}
\end{equation}
Therefore, to insure asymptotic freedom, we consider only
$N_f=0,1,2,3,4$. 

In the previous section, from the counting of fermion zero-modes in an
instanton background, we found that $U(1)_{\cal R}$ is broken to a
discrete subgroup ${\bf Z}_{4N_c-2N_f}$. This was obtained by
requiring the invariance of the correlation function $G$ given by
(\ref{G}) under $U(1)_{\cal R}$. For $N_c=2$, we have an extra
discrete symmetry group (\ref{rho}) which permutes the fermion
zero-modes $\psi_{q_1}$ and $\widetilde\psi_{q_1}$. This group is
anomalous as it changes $G$ to $-G$. Therefore, now we can also allow
$U(1)_{\cal R}$ transformations which do not keep $G$ invariant, but
change it by a sign. This sign can be compensated for by an anomalous
${\bf Z}_2$ transformation. Thus, for $N_c=2$, $U(1)_{\cal R}$ is
broken to the discrete subgroup ${\bf Z}_{2(4N_c-2N_f)}={\bf
Z}_{4(4-N_f)}$. This can be combined with the anomalous ${\bf Z}_2$ to
get a discrete ${\bf Z}_{4(4-N_f)}$ anomaly-free subgroup with the
action (see (\ref{Rb1}) and (\ref{Rb2})):
$$
\begin{array}{l}
\left.\begin{array}{l}
W_{\alpha} \rightarrow \omega W_{\alpha} (\omega^{-1} \theta )\\
\Phi \rightarrow \omega^2 \Phi (\omega^{-1} \theta )\\
Q^1  \rightarrow \widetilde Q_1 (\omega^{-1} \theta )\\
\widetilde Q_1  \rightarrow Q^1 (\omega^{-1} \theta ) \\
\end{array}\right. \\
\left.\begin{array}{l}
Q^i  \rightarrow Q^i (\omega^{-1} \theta )\\
\widetilde Q_i  \rightarrow \widetilde Q_i (\omega^{-1} \theta )\\
\end{array}\right\} \qquad i\not= 1\,,
\end{array}
$$
where, $\omega={\exp}(2i\pi/4(4-N_f))$. For $N_f = 0$ we do not have the
quarks to cancel the anomaly and only the square of the above
transformations is anomaly free. This case was discussed in the
previous section. Furthermore, it can be seen that a subgroup ${\bf
Z}_2 \subset {\bf Z}_{4(4-N_f)}$ is the same as $(-1)^F$ in the
Lorentz group. Combining the above transformations with the $U(1)_J$
subgroup of $SU(2)_R$ (see (\ref{Rb1}) and (\ref{Rb2})), we find a
${\bf Z}_{4(4-N_f)}$ symmetry which commutes $N=1$ supersymmetry
\begin{equation}
\begin{array}{l}
\left.\begin{array}{l}
\Phi \rightarrow \omega^2 \Phi(\theta) \\
Q^1 \rightarrow \omega^{-1} \widetilde{Q}_1(\theta) \\
\widetilde{Q}_1 \rightarrow \omega^{-1} Q^1(\theta) \\
\end{array} \right.\\
\left.\begin{array}{l}
Q^i \rightarrow \omega^{-1} Q^i(\theta) \\
\widetilde{Q}_1 \rightarrow \omega^{-1} \widetilde{Q}_i(\theta) 
\end{array} \right\}
\qquad i \not= 1\,.
\end{array}
\label{symm-A}
\end{equation}
Under this transformation, the gauge invariant order parameter $u={\rm
Tr}\phi^2$ has charge $4$ and transforms as $u \rightarrow \exp 2 \pi
i /(4-N_f) u$. This further breaks ${\bf Z}_{4(4-N_f)}$ down to ${\bf
Z}_4$. The remaining ${\bf Z}_{4-N_f}$ acts non-trivially on the
$u$-plane. Note that for $N_f=0$, the subgroup which does not keep $u$
invariant is only a ${\bf Z}_2$.

As in the $N_f = 0$ case, the large $u$-behaviour of $a_D(u)$ is
determined by the one-loop $\beta$-function (\ref{betanf}):
\begin{eqnarray}
a &\cong& \frac{1}{2} \sqrt{2u} + \cdots \,,\nonumber\\
a_D &\simeq& i \;\frac{4-N_f}{2\pi}\,a(u)\ln\frac{u}{\Lambda^2_{N_f}}
+ \cdots \,.\label{betanf-mon}
\end{eqnarray}
Here, the dots represent non-perturbative instanton corrections. The
generic form of these corrections can be obtained by arguments similar
to the ones used for the $N_f=0$ case in the previous section. First,
a correction coming from a $k$-instanton configuration is proportional
to the $k$-instanton factor, which, using the $\beta$-function
(\ref{betanf}), can be written as 
\begin{equation}
e^{-8\pi^2k/g^2} =\left(\frac{\Lambda_{N_f}}{a}\right)^{k(4-N_f)}\,. 
\label{inst-nf}
\end{equation}
Following Seiberg \cite{Seiberg-n,Seiberg-poh}, we can restore the
broken part of the $U(1)_{\cal R} \times {\bf Z}_2 (\rho )$ symmetry
by assigning appropriate charges to $u$ and $\Lambda_{N_f}$.  Thus we
assign charge $4$ and even $\rho$-parity to $u$ and charge $2 (4-N_f)$
and odd $\rho$-parity to $\Lambda_{N_f}^{4-N_f}$. Note that with this
assignment, the one-instanton factor $\exp(-8\pi^2/g^2)$ will have an
odd $\rho$-parity which compensates for the odd parity under ${\bf
Z}_2(\rho)$ of the correlation function $G$ in (\ref{G}), thus keeping
it invariant . For the special case of $N_f=4$, $U(1)_{\cal R}$ is
non-anomalous and $u$ has charge $4$. $Z_2(\rho)$ in $O(8)$ is still
anomalous, but again, it can be treated as unbroken by assigning odd
parity to the instanton factor $\exp(-8\pi^2/g^2)$. Invariance under
$U(1)_{\cal R}$ with the above charge assignments implies that each
correction term should contain a factor of $\sqrt u$. Moreover, since
the metric on the $u$-plane is invariant under the $\rho$-parity,
configurations with odd instanton numbers cannot contribute to $a$ and
$a_D$. Putting these facts together, we can write the generic form of
the corrected $a$ and $a_D$ as
\begin{eqnarray}
a &=& \frac{1}{2} \sqrt{2u} \left ( 1 + \sum_{n=1}^\infty a_n (N_f)
\left(\frac{\Lambda^2_{N_f}}{u}\right)^{n(4-N_f)} \right )\,,
\nonumber\\
a_D &=& i \frac{4-N_f}{2 \pi} a(u) \ln \frac{u}{\Lambda^2_{N_f}} +
\sqrt{u} \sum_{n=0}^\infty  a_{Dn} (N_f) \left(\frac{\Lambda^2_{N_f}}
{u}\right)^{n(4-N_f)} \,.\nonumber
\end{eqnarray}
The difficult part is to compute the coefficients $a_n$ and $a_{Dn}$.

Since the gauge group $SU(2)$ breaks to $U(1)$, the theory will
have massive charged states and we want to know how the unbroken
global symmetry acts on them. The unbroken part of the global
symmetry, which is the part that keeps $u$ invariant, is obtained by
raising the transformations (\ref{symm-A}) to the power $4-N_f$.
This unbroken transformation changes the sign of $\phi$ and therefore
acts as charge conjugation on the charged fields. For $N_f=1,3$;
$4-N_f$ is odd and the unbroken transformation contains odd powers of
$\rho$. It, therefore, acts as the parity in $O(2N_f)$. Hence, in this
case, the parity (\ref{rho}) is realized on the spectrum but it
reverses the signs of electric and magnetic charges. States of given
charge belong to $SO(2)$ (for $N_f=1$) or $SO(6)$ (for $N_f=3$)
multiplets. For $N_f=2,4$; $4-N_f$ is even and the unbroken
transformation contains only even powers of $\rho$. All the odd
powers of $\rho$, which amount to the parity transformation, are part
of the broken transformations. Thus, in this case, the parity symmetry
is broken and the states are only in $SO(4)$ (for $N_f=2$) or $SO(8)$
(for $N_f=4$) representations. 

\subsection{BPS Saturated States}

On the Higgs branch $SU(2)$, is completely broken and there are no
electric or magnetic charges. Thus the central charge of the $N=2$
algebra only contains contributions from the $U(1)$ charges of the
hypermultiplets. We will not discuss this in any detail. 

On the Coulomb branch, the simplest BPS saturated states are the
quarks with zero bare mass and which acquire masses $M = \sqrt{2} |a|$
after the spontaneous breaking of the gauge symmetry. These form a set
of BPS states which transform as a vector of $SO(2N_f)$. Besides
these, there are BPS states which transform as a spinor of $SO(2N_f)$
arising as follows: Since the gauge symmetry is broken to $U(1)$, the
theory contains monopoles. In the presence of a monopole each $SU(2)$
doublet of fermions has one zero-mode. Since there are $N_f$
hypermultiplets, there are $2N_f$ fermion doublets and, therefore,
$2N_f$ fermion zero-modes. After quantization, these zero-modes give
rise to a $2N_f$-dimensional Dirac algebra which provides a spinor
representation of $SO(2N_f)$. Thus the fermion zero-modes turn the
monopole into a spinor of $SO(2N_f)$. This is very similar to the
quantization of the Ramond sector of the Superstring theory. The
presence of spinors indicate that, at the quantum level, the symmetry
group is a universal cover of $SO(2N_f)$, or $Spin(2N_f)$.

One of the monopole's collective coordinates is a charge rotation.
Upon quantization this leads to a spectrum of electrically charged
states for the monopole. A $2\pi$ rotation by the electric charge
operator, however, is not the identity.  Using the Witten effect, such
a rotation gives a topologically non-trivial gauge transformation with
eigenvalue $e^{i \theta} (-1)^H$ for a monopole with $n_m = 1$. Here,
$(-1)^H$ is the centre of $SU(2)$ which is odd for states in the
hypermultiplet and even for states in the vector multipler. In this
section we normalize the charge operator so that elementary quarks
have charges $\pm 1$ and massive gauge bosons have charges $\pm 2$. 
With this normalization, the above mentioned gauge rotation can be
written as an operator statement 
$$
e^{i \pi Q} = e^{i \theta n_m} (-1)^H.
$$
where, $Q = n_e + n_m \theta/\pi$. Since, $n_e \in {\bf Z}$, this
relation implies that states of even $n_e$ have $(-1)^H = 1$ and
states of odd $n_e$ have $(-1)^H = -1$. The factor $(-1)^H$ can be
identified with the chirality operator for the spinor representation
of $SO(2N_f)$.  For $N_f = 1,3$, the $SO(2N_f)$ parity transformation 
acts on the spectrum and it guarantees that a dyon transforming as a
positive chirality spinor of $SO(2N_f)$ is degenerate with a particle
of opposite electric and magnetic charges and opposite $SO(2N_f)$
chirality.  No such relation exists for $N_f = 2,4$.

For $N_f > 0$ we will see that the spectrum contains states with $n_m
> 1$. A convenient way of labelling states is in terms of their
behaviour under the centre of $Spin(2N_f)$.
\begin{itemize}
\item

$N_f = 2$: \, $Spin(4) = SU(2) \times SU(2)$ with centre ${\bf Z}_2
\times {\bf Z}_2$. The representations of the center can be labeled
by $(\epsilon , \epsilon ')$, where $\epsilon = 0$ for vector-like
irreducible representation and $\epsilon = 1$ for spinor-like
irreducible representations.  An elementary quark transforms as
$(\epsilon , \epsilon ') = (1,1)$. Multiple quark states then
transform as $(n_e \; mod \; 2, n_e \; mod \; 2)$. Since monopoles
behave like spinors, we have
$$
(\epsilon , \epsilon ') = ((n_e + n_m) \; mod \; 2, n_e \; mod \; 2)
\qquad {\rm for}\,\, {N_f = 2}\,.
$$
\item
$N_f = 3$: \, $Spin(6) = SU(4)$ with centre ${\bf Z}_{4}$.  The
elementary quarks are in the $6$ of $Spin (6)$, thus they have charge
$2$ with respect to the centre.  Since monopoles behave like spinors,
we conclude that ${\bf Z}_{4}$ acts as:
$$
\exp \; \frac{i \pi}{4} (n_m + 2n_e)
\qquad{\rm for} \,\,{N_f = 3}\,.
$$
\item
$N_f=4$:\, $Spin(8)$ has centre ${\bf Z}_2\times {\bf Z}_2$. Hence,
using similar arguments as above, we get,
$(\epsilon,\epsilon')$$=$$(n_m\; mod\; 2, n_e\; mod \; 2) $. The four
representations of the center are labeled by the representations of of
$Spin(8)$ which realize them: $(0,1)\equiv v$, $(1,0)\equiv s$,
$(1,1)\equiv c$ and $(0,0)\equiv o$.  Here, $v$ refers to the vector
representation, $s$ and $c$ to two spinor representatios and $o$ to
the trivial representation. $Spin(8)$ has a triality group of outer
automorphisms which is isomorphic to the permutation group ${\bf S}_3$
acting on $v$, $s$ and $c$.
\end{itemize}

If an $N = 2$ invariant mass is turned on, say $m_{N_f} \not= 0$,
then $SO(2N_f)$ explicitly breaks to $SO(2N_f-2) \times SO(2)$, and
the global Abelian charge associated with $SO(2)$ appears in the
central charge 
$$
Z = n_e a + n_m a_D +S_{N_f} \frac{m_{N_f}}{\sqrt{2}},\qquad
M = \sqrt{2} |Z|\,.
$$
Hence for $a = \pm m_{N_f} / \sqrt{2}$ one of the elementary quarks
becomes massless. This will be crucial later in determining the global
structure of the quantum moduli space on the Coulomb branch.

\subsection{Duality}

As in the $N_f=0$ case, there is an $SL(2,Z)$ duality group which acts
on the fields and the couplings. In the presence of matter, some new
issues appear which will be discussed below.

As before, we can compute the monodromy matrices for $(a_D,a)$ around
the singular points on the moduli space.  The simplest case
corresponds to having a single quark with non-zero bare mass, and
investigating what happens as $a \approx a_0 \equiv m_{N_f} /
\sqrt{2}$.  From the QED analysis we know that there is a logarithmic
singularity at $a_0$ where this quark becomes massless. Near the
singularity we have 
$$
\begin{array}{rcl}
a& \approx & a_0\,, \\
a_D & \approx & - \frac{i}{2\pi} (a-a_0) \ln (a-a_0) +
c\,. 
\end{array}
$$
The monodromy can now be easily computed to be
$$
\begin{array}{rcl}
a &\rightarrow & a \,,\\
a_D &\rightarrow &\displaystyle{ a_D + a - \frac{m_{N_f}}{\sqrt{2}}}\,.
\end{array}
$$
Unlike the $N_f = 0$ case, now we have an inhomogeneous transformation
as the column vector $(a_D, a)^T$ picks up a shift under the monodromy
(besides the usual $SL(2,Z)$ transformation). The possibility of such
a shift was also noticed for the pure gauge theory case in the
previous section. There, however, the shift was not part of the
monodromy group and, furthermore, it was not compatible with the
symmetries of the BPS mass formula. In the presence of matter, this
shift is allowed since the BPS mass formula is modified. Moreover, now
the shift naturally appears as a part of the monodromy group. To write
a monodromy matrix, we construct a column vector $(m/ \sqrt{2} , a_D ,
a )^T$. The monodromy can now be written as 
$$
\left ( \begin{array}{c} m/ \sqrt{2} \\ a_D \\ a \end{array} \right)
\rightarrow {\cal M} 
\left ( \begin{array}{c} m/ \sqrt{2} \\ a_D \\ a \end{array}\right)\,, 
$$
with the monodromy matrix ${\cal M}$ given by 
$$
{\cal M } = \left ( \begin{array}{ccc}
1 & 0 & 0 \\ -1 & 1 & 1 \\ 0 & 0 & 1
\end{array} \right )\,,\qquad
{\cal M}^{-1} = \left(\begin{array}{ccc}
1 & 0 & 0 \\ 1 & 1 & -1 \\ 0 & 1 & 1
\end{array}\right) \,.
$$
If we arrange the charges as a row vector $W = (S, n_m, n_e)$, then
invariance of $Z$ (or of $M$) implies that, under the monodromy,
$W \rightarrow W{\cal M}^{-1}$. In general the matrix ${\cal M}$ can
be of the form
$$
{\cal M} = \left ( \begin{array}{ccc}
1 & 0 & 0 \\ r & k & l \\ q & n & p
\end{array} \right )\,,
\qquad
{\cal M}^{-1} = \left ( \begin{array}{ccc}
1 & 0 & 0 \\ lq-pr & p & -l \\
nr-kq & -n & k
\end{array} \right )\,,
$$
with $\det {\cal M} = 1$. 

Note that under the monodromy, the electric and magnetic charges can
mix with each other but will not pick up contributions proportional to
$S$ which is a global symmetry charge.  On the other hand, $S$ can
pick up contributions proportional to $n_e$ and $n_m$ which are
related to the local gauge symmetry.  For the example considered
above,
$$
(S, n_m, n_e ) \rightarrow (S + n_m, n_m, -n_m + n_e )\,.
$$
This leads once again to issues of marginal stability: For large
values of $m_{N_f}$, the singularity is in the weak-coupling region of
large $u$ where a semi-classical treatment of monopole is reliable.
In the semi-classical quantization, the global $U(1)$ charge is
carried only by fermionic zero-modes. Since there is only a finite
number of these modes, one cannot construct states with arbitrarily
large values of $S$. This means that, although, by acting with ${\cal
M}^{-1}$ we can increase the value of $S$ at will, what might have
started as a one-particle state comes back as a multiparticle state.  
This is possible if at some point, when going around the singularity,
the single-particle state becomes unstable and decays into a
multi-particle state. Thus, although the formalism is $SL(2, Z)$
covariant, the spectrum is not. 

Seiberg and Witten provide many compelling arguments to suggest that
the theory with $N_f = 4$, like the theory with $N = 4$, is
$SL(2,Z)$ invariant, though the details are quite
different. In the $N_f=4$ theory, the global symmetry group is
$Spin(8)$ which is the universal cover of $SO(8)$. The states $(n_m,
n_e)=(0,1)$ are the elementary hypermultiplets in the vector
representation $v$ of $Spin(8)$.  The states $(n_m, n_e)=(1,0)$ are in
a spinor representaion $s$, and $(n_m, n_e)=(1,1)$ are in a spinor
representation $c$ of $Spin(8)$. While $SL(2,Z)$ alone cannot keep
this spectrum unchanged, a combination of $SL(2,Z)$ and the $Spin(8)$
triality group (which permutes the representations $v$, $s$ and $c$)
could do so. This is provided one is willing to assume that 
there are monopole bound states for every pair of relatively prime
integers $(p,q)$, analoguous to the situation discussed by Sen for the
$N=4$ theory \cite{sen}.  For each such pair, there should exist eight
states transforming in a representation $8_v$, $8_s$ or $8_c$ of
$Spin(8)$, depending on the $mod \; 2$ reduction of $(p,q)$.
$SL(2,Z)$ mixed with triality will then keep the spectrum
invariant. The solution Seiberg and Witten provided for this model
gives strong support to this possibility.  The global symmetry group
is then a semi-direct product of $Spin(8)$ and $SL(2,Z)$.

\subsection{A First Look at Singularities}

As in the pure $N = 2$, we first try to locate the singularities on
the moduli space before computing their monodromies.  We recall that,
with the standard normalization of the gauge coupling constant $g$,
the one-loop $\beta$-function which is given by (\ref{betanf}), can be
integrated to give
$$
\frac{1}{\alpha_{N_f}(\mu)} = \frac{4-N_f}{2 \pi} \ln 
\frac{\mu}{\Lambda_{N_f}}\,,
$$
where, $\alpha=4\pi/g^2$.  Now, if some (say $N_f-N_f'$) of the quarks
have masses $m_i=m$, and we are looking at the theory at some scale
$\mu < m$, then the low energy theory contains only $N_f'$
hypermultiplets as the degrees of freedom. The coupling $\alpha(N_f')$
is then given by an expression similar to the above one with $N_f$ and
$\Lambda_{N_f}$ replaced by $N_f'$ and $\Lambda_{N_f'}$. The scales of
the theories can be related by the matching condition $\alpha_{N_f}(m)
=\alpha_{N_f'}(m)$, which implies
$$
\Lambda^{4-N_f'}_{N_f'} = m^{N_f - N_f'}\;\Lambda_{N_f}^{4-N_f}\,.
$$
For instance,
$$
\begin{array}{lcl}
N_f = 3 \;, N'_f = 0\, &\Rightarrow& \Lambda^{4}_{0} = m^3 \Lambda_3\,,\\
N_f = 1 \;, N'_f = 0\, &\Rightarrow& \Lambda^{4}_{0} = m \Lambda_{1}^{3}\,.
\end{array}
$$

To determine the singularity structure on the moduli space, we first
consider theories with $N_f \leq 3$, and with hypermultiplet bare
masses very large as compared to $\Lambda$. The singularities which
arise from hypermultiplets becoming massless are now in the
semi-classical (large $u$) region of the moduli space and can be easily
identified. In the small $u$ region, one is effectively left with an
$N_f=0$ theory with two singularities corresponding to massless
monopoles and dyons. Then we slowly decrease the bare masses to zero
and follow the movement of the sigularities on the moduli space.

\noindent\underline{$N_f = 3$}:
Let us start with equal masses $m_i = m >> \Lambda$. In this case, the
global symmetry $Spin(6)\approx SU(4)$ of the massless theory is
broken to $SU(3) \times U(1)$. Classically, there is a singularity at
$a = m/ \sqrt{2}$ where the three quarks become massless. These
electrically charged massless fields form a \underline{3}
representation of $SU(3)$. For $m >> \Lambda$ the singularity is in
the semiclassical region, $u \sim 2a^2 >> \Lambda^2$. For $u<< m^2$
the three massive quarks can be integrated out giving a $N_f=0$ theory
with scale $\Lambda^{4}_{0} = m^3 \Lambda_3$. Hence, for small $u$ the
moduli space is that of a pure $N = 2$ theory with scale $\Lambda_{0}$
which has two singularities with $(n_m , n_e) = (1,0)$ and
$(1,1)$. These are the points where monopoles and dyons become
massless. Clearly the massless states at these singularities are
$SU(3)$ invariant. 

As $m$ is decreased, the singularity at large $u$ moves, and although
the charges of the associated states can change (through monodromy
matrices), their non-Abelian charges cannot change. Hence, for any
$m$, the massless fields at the various singularities transform as
\underline{3}, \underline{1} and \underline{1} of $SU(3)$. In the $m =
0$ limit, the original global symmetry is restored and the massless
states must combine into representations of $SU(4)$. The only
possibility is to have two singularities combining into a
\underline{4} of $SU(4)$ and the other sigularity moves somewhere else
remaining a singlet. Thus, we conclude that the massless $N_f = 3$
theory has two singularities with massless particles in the
\underline{4} and \underline{1} of $SU(4)$. From our study of how
different states transform under the centre, the smallest charges for
the massless particles at the singularities are $(n_m , n_e) = (1,0)$
for the \underline{4}, and $(n_m, n_e) = (2,1)$ for the \underline{1}.
Semiclassically the first state exists, but it is not known whether
the monopole-monopole bound state implied by the second also exists or
not.

\noindent\underline{$N_f = 2$}:
In this case, with two equal masses $m_i=m >>\Lambda$, the $Spin(4)$
global symmetry of the massless theory is broken to $SO(2)
\times SO(2)$. There is a singularity at $a = m/ \sqrt{2}$ and the
massless states there transform under one or the other of the two
$SO(2)$'s. In the region $u<<m^2$, we can again integrate out the
quarks, leading to a $N_f=0$ theory with $\Lambda^4_0=m^2\Lambda^2_2$,
and two singularities with $(n_m, n_e) = (1,0), (1,1)$. The massless
fields at these singularities are singlets under $SO(2)\times SO(2)$.
In all, we have four massless states associated with three
singularities. As $m \rightarrow 0$, we recover the full global
symmetry group of the massless theory with two flavours which is 
$Spin (4) \approx SU(2) \times SU(2)$. The singularities have to
combine in such a way that the massless states form a representation
of this unbroken group. Since each $SO(2)$ is contained in a different
$SU(2)$, the only way to combine the massless states into
representations of $SU(2)\times SU(2)$ is as (\underline{2},
\underline{1}) and (\underline{1}, \underline{2}). Hence, for $m = 0$,
there are two singularities, and the simplest charge assignments are
$(n_m , n_e ) = (1,0)$ in one spinor of $SO(4)$, and $(n_m , n_e ) =
(1,1)$ in the other spinor. Recall that the transformation under the
centre is $(\epsilon,\epsilon')=((n_m+n_e) mod\; 2, n_e\; mod\;2)$. 

\noindent\underline{$N_f=1$}:
Now the massive theory has the same $SO(2)$ symmetry as the massless
theory, however, the same arguments as before imply the presence of
three singularities for large $m$. This number does not change as
$m\rightarrow 0$ due to the ${\bf Z}_3$ symmetry acting on the
$u$-plane. We recall that this symmetry is the subgroup ${\bf
Z}_{4-N_f}={\bf Z}_{3}$ of ${\bf Z}_{4(4-N_f)}$ in (\ref{symm-A})
which is broken by a non-zero $u$. From 
$$
a_D = i \frac{4-N_f}{2\pi} a(u) \ln \frac{u}{\Lambda^2_{N_f}}+...\,,
$$
we see that as $a$ transforms homogeneously under (\ref{symm-A}),
$a_D$ is shifted by $a$, {\it i.e.}, $a\rightarrow \omega^2 a$, $a_D
\rightarrow \omega^2 (a_D + a)$ with $\omega = e^{i\pi/6}$. Therefore,
if one of the singularities in the $m=0$ limit corresponds to massless
states with $(n_m , n_e ) = (1,0)$, then the ${\bf Z}_3$ symmetry
implies the existance of two other singularities characterized by
$(n_m , n_e )= (1,1)$ and $(1,2)$. Hence, even for $m=0$, the moduli
space of the $N_f=1$ theory has three singularities corresponding to
massless states $(n_m , n_e )= (1,0), (1,1)$ and $(1,2)$.

\subsection{Monodromies and the Determination of the Metric}

As in the case of $N=2$ theory without matter, a solution to the
theory can be found by regarding the $u$-plane as the moduli space of
a family $E_u$ of elliptic curves parametrized by $u$. The quantities
$a_D(u)$ and $a(u)$ can then be related to the periods of these
curves. A curve in $E_u$ becomes singular when one of the cycles on it
shrinks to zero size. This corresponds to a singularity on the
$u$-plane with a non-trivial monodromy around it. If we know the
singularities and the associated monodromies on the $u$-plane (which
are associated with the appearence of massless particles in the
spectrum), then we can work backwords and determine the families of the
elliptic curves from which the periods $a_D$ and $a$ could be
computed. In this subsection, we sketch the physical arguments
used by Seiberg and Witten in \cite{SW2} to find the curves for
$N_f=1,2,3$ theories exploiting the general features of the curve for
the $N_f=0$ theory (This reference also contains a more rigorous
treatment of this problem which we will not reproduce here).

For $N_f = 0$, the solution was given by a family of elliptic curves
described by the equation
$$
y^2 = (x - \Lambda^2 )(x + \Lambda ^2 )(x-u) \,,
$$
with the monodromy in $\Gamma(2)$. In the present section we have
changed our conventions so that $(n_m , n_e )$ are both integers even
in the presence of matter fields. This was achieved by multiplying
$n_e$ by $2$ and dividing $a$ by $2$ so that $Z = a_D n_m + a n_e$ is
unchanged. This change of convention can be implemented as a
transformation  
$$
\left ( \begin{array}{cc} a_D \\ a \end{array} \right ) \rightarrow
\left ( \begin{array}{cc} a_D \\ a' \end{array} \right ) =
\left ( \begin{array}{cc} 1 & 0 \\ 0 & 1/2 \end{array} \right )
\left ( \begin{array}{cc} a_D \\ a \end{array} \right ),
$$
which also changes the monodromy matrix as
$$
\left ( \begin{array}{cc} m & n \\ p & q \end{array} \right )
\rightarrow
\left ( \begin{array}{cc} 1 & 0 \\ 0 & 1/2 \end{array} \right )
\left ( \begin{array}{cc} m & n \\ p & q \end{array} \right )
\left ( \begin{array}{cc} 1 & 0 \\ 0 & 2 \end{array} \right ) =
\left ( \begin{array}{cc} m & 2n \\ p/2 & q \end{array} \right )\,.
$$
Since $p$ and $n$ are integers modulo 2, the new monodromy matrix
contains the entry $2n$ which is an integer modulo 4, while all other
entries are integers. These martices form the subgroup $\Gamma_{0}(4)$
of $SL(2,Z)$. In the new convention, the monodromies in (\ref{mondall})
take the form
\begin{equation}
M_{\infty} = \left( \begin{array}{cc}
-1 & 4 \\ 0 & -1 \end{array} \right )\,, \quad
M_{1} = \left( \begin{array}{cc}
1 & 0 \\ -1 & 1 \end{array} \right )\,, \quad
M_{-1} = \left( \begin{array}{cc}
-1 & 4 \\ -1 & 3 \end{array} \right )\,,
\end{equation}
and the solution is given by a family of curves described by  
\begin{equation}
y^2 = x^3 - ux^2 + \frac{1}{4} \Lambda^4 x \,.
\label{f2}
\end{equation}
As discussed in the previous section, any genus-one curve can be
represented by a cubic
$$
y^2 = F(x) = (x-e_1)(x-e_2) (x-e_3)\,.
$$
This describes a space $x$ as a double cover of the complex plane
branched over $e_1$, $e_2$, $e_3$ and $\infty$. The curve becomes
singular when two branch points coincide ({\it e.g.}, $e_1=e_2\not=
e_3$ or $e_3 \rightarrow \infty$). In this case the singularity is
called stable.  If more than two branch points coincide, then the
singularity is not stable, but a $u$-dependent reparametrization of
$x$ and $y$ can always convert this into a stable singularity. For the
curve (\ref{f2}) the branch points are at 
$$
x=0,\, \frac{1}{2} (u \pm \sqrt{u^2 - \Lambda^4}),\, \infty\,.
$$
For $u = \pm \Lambda^2$ we have two stable singularities, but the $u
\rightarrow \infty$ singularity is not stable.  

To understand the properties of this curve better, let us first
consider a generic situation: For a stable singularity at, for
instance, $u = 0$, the family of curves near $u=0$ can be written in
the form 
$$
y^2 = (x-1)(x^2 - u^n)\,, 
$$
for some integer $n$. The monodromy around $u=0$ is then conjugate to
$T^n$ where the matrix $T$ is defined in (\ref{ST}). This can be
understood as follows: Consider the holomorphic Abelian differential
$\omega = dx/y$ on a curve $y^2=(x-1)(x^2-\lambda)$, where $\lambda =
u^n$. The periods can be written as  
$$
\omega_1 = \int_{\Gamma_1}\;\frac{dx}{y}\,, \qquad
\omega_2 = \int_{\Gamma_2}\;\frac{dx}{y}\,,
$$
where $\Gamma_1$ is a path from $u=-\lambda^{1/2}$ to
$u=\lambda^{1/2}$, and $\Gamma_2$ is a path from $u=\lambda^{1/2}$ to  
$u=1$ (Fig. 6(a)). As $\lambda \rightarrow e^{2 \pi i} \lambda$,
$\lambda^{1/2} \rightarrow -\lambda^{1/2}$ and the cut $\Gamma_1$ 
moves as in Fig. 6(b). This simply exchanges the branches of the
integrand and therefore $\omega_1 \rightarrow \omega_1$.  


\noindent As for $\omega_2$, the path $\Gamma_2$ is transformed as in
Fig. 7.  

\noindent Thus, $\omega_2 \rightarrow \omega_2 + \omega_1$. Since
$\lambda = u^n$, when $u \rightarrow e^{2\pi i}u$, $\lambda$ makes $n$
turns and we obtain the $n$-th power of the monodromy. Thus, in terms
of $\lambda$, the monodromy is conjugate to $T$ while in terms of $u$,
it is conjugate to $T^n$ as we wanted to show. 


% ==== FIGURE_SLOT_6 ====

% ==== FIGURE_SLOT_7 ====


Given any polynomial, a useful quantity is its discriminant defined as
$$
\Delta = \prod_{i<j} (e_i - e_j)^2,
$$
where the $e_i$'s are the roots.  $\Delta$ can be expressed in terms of
the coefficients of the polynomial. Clearly, at a singularity when two
branch points coincide, $\Delta = 0$ (except for a singularity at
$\infty$). For the example $y^2 = (x-1)(x^2-u^n)$, $\Delta \sim u^n$
near $u = 0$ where the monodromy is conjugate to $T^n$. Hence, in the
stable case, the exponent of the monodromy is the order of zero of
$\Delta$. For instance, for the family of curves described by
(\ref{f2}), $\Delta$ has first-order zeros at $u=\pm\Lambda^2$, and
the monodromies around these points are conjugate to $T$. 
Let us now look at $u\rightarrow \infty$. For large $u$ the branch
points are approximately at $x=0,\,\Lambda^4/4u,\,u,\,\infty$. Thus
the singularity at $u=\infty$ is not stable as more than two branch
points coincide in this limit. This is cured by a change of variables 
$$
x = x'u\,, \qquad y = y' u^{3/2}\,,
$$
which shifts the branch points to $0,\,\Lambda^4/4u^2,\,1,\,\infty$.
The discriminant now behaves like $\Delta \sim u^{-4}$ corresponding
to a monodromy conjugate to $T^{-4}$. Due to the presence of $\sqrt u$
in $y \rightarrow y'$, the monodromy in terms of the original
varialbles $(x,y)$ has an extra minus sign and is conjugate to
$PT^{-4}$. Thus we see that the curve (\ref{f2}) produces the correct
monodromies of the $N_f = 0$ theory.

To obtain the curves for theories with non-zero $N_f$, we should know
the monodromies which arise in these theories. This can be easily
worked out (as in the $N_f=0$ case) since we already know the charge
spectrum of the particles which become massless at the singularities
on the $u$-plane. First note that, as in the $N_f=0$ theory, the
monodromy at $u=\infty$ can be obtained form the perturbative
$\beta$-function (\ref{betanf}), or equivalently (\ref{betanf-mon}),
and is given by 
$$
{\cal M}_{\infty} = PT^{N_f -4}.
$$
The singularities at finite points in the $u$-plane, in general,
correspond to massless magnetic monopoles with one unit of magnetic
charge and $n_e$ units of magnetic charge. To calculate the
corresponding monodromy, we have to go to a dual description of the
theory in which the monopole couples to the dual gauge field the way
an electron couples to the usual gauge field. 
In this frame, the monodromy is determined by the one-loop QED 
$\beta$-function (\ref{beta-hyper}) with $k$ hypermultiplets and is
given by $T^k$. This has to be conjugated with $T^{n_e}S$ which
converts a hypermultiplet of charge $(0,1)$ into a monopole of charge
$(1,n_e)$. Hence the monodromy around a point with $k$ magnetic
monopoles with $(n_m = 1, n_e)$ is $(T^{n_e}S) T^k (T^{n_e}S)^{-1}$.  
Similar arguments can be applied to calculate the monodromy for the
$(2,1)$ state in the $N_f = 3$ theory. In the following we list
all the monodromies around the singularities described in the previous
subsection: 
$$     
\begin{array}{rl}
N_f=0:& STS^{-1},\,(T^2 S)T(T^2 S)^{-1}\rightarrow M_{\infty}=PT^{-4}\,,\\
N_f=1:& STS^{-1},\,(TS)T(TS)^{-1},\, (T^2S)T(T^2S)^{-1} \rightarrow
M_{\infty}=PT^{-3}\,, \\
N_f=2:& ST^2 S^{-1},\,(TS)T^2(TS)^{-1}\rightarrow M_{\infty}=PT^{-2}\,,\\
N_f=3:& (ST^2S)T(ST^2S)^{-1},\,ST^4 S^{-1}\rightarrow
M_{\infty}=PT^{-1}\,. \\
\end{array}
$$
Using $S^2=-1$ and $(ST)^3=1$, it is easy to check that, for each
$N_f$, the product of the monodromy matrices at finite $u$ yields
$M_{\infty}$.  Note that in the first three cases ($N_f =
0,1,2$) the electric charges assigned to the singularities differ in
sign from those determined by our previous arguments.  This is
consistent with the inherent ambiguities in relating the universal
charges with the transformation properties under the centre of the
group. Based on the knowledge obtained so far, in the following we
sketch the arguments of Seiberg and Witten to obtain the explicit form
of the curves for theories with non-zero $N_f$.

\noindent\underline{Properties of the $N_f =0$ Curve:}
Before proceeding further, it is very useful to enumerate some
properties of the $N_f=0$ curve which are expected to remain valid
even for non-zero $N_f$: 
\begin{enumerate}
\item
The equation describing the family of curves is of the generic form
$y^2=F(x,u,\Lambda)$, where $F$ is a polynomial at most cubic in $x$
and $u$.
\item
The part of $F$ cubic in $x$ and $u$ is $F_0 = x^2 (x-u)$.
\item
If we assign $U(1)_{\cal R}$ charge $4$ to $u$ and $x$ and charge $2$
to $\Lambda$, then $F$ has charge $12$. If $y$ is assigned charge $6$,
then the curve is invariant under $U(1)_{\cal R}$ transformations. 
\item
$F$ can be written as $F = F_{0} + \Lambda^4 F_1$ where $F_1 = x/4$.
\end{enumerate}

The property (1) remains an ansatz for $N_f > 0$. It will lead to
the correct monodromy matrices, and it can also be justified in part
when considering the $N_f = 4$ theory.  Property (2) is a consequence
of the fact that as $u \rightarrow \infty$, we must obtain the
monodromy at infinity associated with a logarithm in $a_D$ coming
from the one-loop $\beta$-function. This means that one of the branch
points at finite $x$ should move to infinity as $u\rightarrow\infty$. 
For large $u$, the cubic part can be written as $(x-e_1 u)(x-e_2
u)(x-e_3 u)$. Then the desired behaviour can be obtained if two
$e_i$'s coincide and the other is different, say $e_1=e_2\not=e_3$.
In this case, by a redefinition of $x$, we can bring $F_0$ to the
form $x^2 (x-u)$. Property (4) in the $N_f = 0$ theory is a
consequence of the $U(1)_R$ charge assignments.  Note that $F$ has
only a classical contribution plus a one-instanton term : $\Lambda^4
x/4$.  Now we use these properties to determine the curves for 
$N_f = 1,2,3$.

\noindent\underline{The Curves for Massless $N_f=1$ Theory:}
In this theory, from equation (\ref{inst-nf}), the instanton amplitude
is proportional to $\Lambda^3_{1}$. However, for $N_f \geq 1$, the
instanton factor is odd under the $\rho$ parity and, therefore, only
even powers of it can appear in $F$. Since $\Lambda_{1}$ has $U(1)_R$
charge $2$ and $F$ has charge $12$, the only possibility is 
$$
y^2 = x^2 (x-u) + t \Lambda^{6}_{1} \,.
$$
The constant $t$ can be absorbed in a redefinition of $\Lambda_{1}$:
$t\Lambda^{6}_{1}=\tilde{\Lambda}^6$. The discriminant of this family
of curves is
$$
\Delta = \tilde{\Lambda}^{6}_{1} (4 u^3 - 27 \tilde{\Lambda}^{6}_{1})\,.
$$
This has three zeros which are interchanged under the ${\bf Z}_3$
transformation acting on the $u$-plane and the associated monodromies
are conjugate to $T$. Similarly, the monodromy at large $u$ can be
worked out to be $PT^{-3}$. This is consistent with what we should
have for the $N_f=1$ theory.

\noindent\underline{The Curves for Massless $N_f=2$ Theory:}
In this case the instanton factor is $\Lambda^{2}_{2}$ and, again, in
the absence of bare masses, only even powers of it can appear in $F$. 
Since $\Lambda_2$ has $U(1)_R$ charge 2, the general form of the curve
is
$$
y^2 = x^2 (x-u) + (ax + bu) \Lambda^{4}_{2}\,.
$$
From our discussion of the $N_f=2$ theory, we expect two singularities,
each with two massless monopoles.  Hence the monodromy at each
singularity is conjugate to $T^2$ which means that the discriminant at
each finite singularity should have a double zero.  This condition
determines $a$ and $b$.  After a rescaling of $\Lambda_2$, the
family of $N_f=2$ curves can be written as
$$
y^2 = (x^2 - \tilde{\Lambda}^{4}_{2} ) (x-u)\,,
$$
which has the expected ${\bf Z}_2$ symmetry.

\noindent\underline{The curves for Massless $N_f=3$ Theory:} In this
case, there are two singularities on the $u$-plane with monodromies
conjugate to $T^4$ and $T$ respectively, and there is no symmetry
acting on the $u$-plane. Take the $T^4$ singularity to be at $u =
0$. The discriminant then should have a fourth-order zero at
$u=0$. This, together with the usual $U(1)_{\cal R}$ charge and $\rho$
parity assignments, leads to
$$
F = a \Lambda^{2}_{3} x^2 + bu^2 x + cu x^2 + x^3 \,.
$$
Here, $b \not= 0$ since otherwise the curve is singular for all $u$.
Requiring that the cubic part of $F$ have the expected classical
behaviour, and after some rescaling, one gets
$$
y^2 = x^2 (x-u) + \tilde {\Lambda}^{2}_{3} (x-u)^2\,.
$$

Note that in the above we have only considered theories without a bare
mass term. Since a mass term has odd $\rho$-parity, in a theory with
non-zero bare masses, odd powers of the instanton factor can also
contribute to the equation for the curve. We will not discuss these
cases here, but for the sake of completeness, will simply reproduce
the final results:
\begin{eqnarray}
N_f=1:\qquad y^2 &=& x^2(x-u)+ \frac{1}{4}m\Lambda_1^3 x -
\frac{1}{64} \Lambda_1^6\,, \nonumber \\
N_f=2:\qquad y^2 &=& (x^2-\frac{1}{64}\Lambda_2^4) (x-u) + m_1 m_2
\Lambda_2^2 x - \frac{1}{64}(m_1^2+m_2^2)\Lambda_2^4\,,\nonumber \\
N_f=3:\qquad y^2 &=& x^2(x-u)-\frac{1}{64} \Lambda_3^2(x-u)^2
-\frac{1}{64}(m_1^2+m_2^2+m_3^2) \Lambda_3^2(x-u) \nonumber \\
      && +\frac{1}{4}m_1m_2m_3 \Lambda_3 x - \frac{1}{64}
(m_1^2 m_2^2 + m_2^2 m_3^2 + m_1^2 m_3^2)\Lambda_3^2\,. \nonumber
\end{eqnarray}
For more details and discussions (including the $N_f=4$ case), as well
as a more rigorous method of obtaining the equations for the curves,
the reader is referred to the original work of Seiberg and Witten
\cite{SW2}.   

Once the curves are determined, the rest of the procedure is very
similar to the $N_f=0$ case. We define the quantities $a$ and $a_D$ by
the contour integrals
$$
a=\int_{\gamma_1} \, \lambda\,,\quad a_D=\int_{\gamma_2}\lambda\,,
$$
and the metric is given by ${\rm Im}(\tau)$ with $\tau=\frac{da_D}{du}/
\frac{da}{du}$. Here $\lambda$ is a holomorphic one-form such that
$$ 
\frac{d\lambda}{du}=\frac{\sqrt 2}{8\pi}\frac{dx}{y}\,.
$$
The difference with the $N_f=0$ case is that now $\lambda$ can have
poles with non-zero residue, however, the residues should be $u$
independent. If this is the case, then $\tau$ still undergoes an
$SL(2,Z)$ transformation while $(a_D, a)$ undergoes an $SL(2,Z)$
transformation plus a shift. As we saw in the case of QED with quarks
of non-zero bare masses, this shift is actually needed and is
consistent with the modified form of the BPS mass formula
(\ref{bps-m}). From this it follows that the residues of $\lambda$,
which are responsible for the shifts, should be proportional to the
bare quark masses. This was checked in \cite{SW2} for the $N_f=2$
curve and found to be the case. In fact, the existence of residues is
very restrictive and this information was used by Seiberg and Witten
to drive the curve for the $N_f=4$ theory in a rigorous way. The
curves for other theories can then be determined by renormalization
group flow and the results agree with what is listed above.
